% === Begin Label Data ===
% ID: 791
% 难度星级: 
% 标签: 
% 备注: 
% 组卷引用次数: 0
% === End  Label Data ===

\begin{problem}{2019}{G}{上海卷}{16}{三角函数}
已知 $\tan\alpha\cdot\tan\beta=\tan(\alpha+\beta)$．给出下面两个命题：  
\circled{1}存在角 $\alpha$ 在第一象限，角 $\beta$ 在第三象限；  
\circled{2}存在角 $\alpha$ 在第二象限，角 $\beta$ 在第四象限；  
则 (\hspace{1cm})
\begin{choices}
\choice{{\circled{1}\circled{2}均正确}}
\choice{{\circled{1}\circled{2}均错误}}
\choice{{\circled{1}正确，\circled{2}错误}}
\choice{{\circled{1}错误，\circled{2}正确}}
\end{choices}
\end{problem}

\begin{answer}
D
\end{answer}

\begin{solutions}
 由条件，$\tan\alpha\cdot\tan\beta=\tan(\alpha+\beta)=\dfrac{\tan\alpha+\tan\beta}{1-\tan\alpha\tan\beta}$．  
记 $t=\tan\alpha\tan\beta$，则 $\tan\alpha+\tan\beta=t(1-t)$．  

(1) 若 $\alpha$、$\beta$ 分别在第一、第三象限，则 $t>0$ 且 $\tan\alpha+\tan\beta>0$，从而  
$$
t(1-t)=\tan\alpha+\tan\beta\geq 2\sqrt{\tan\alpha\tan\beta}=2\sqrt{t}\geq 0.
$$  
所以 $0<t\leq 1$，此时 $0\leq t(1-t)-2\sqrt{t}=(\sqrt{t}-1)^2-t^2-1<1-t^2-1=-t^2<0$，矛盾．所以 $\alpha$、$\beta$ 不可能分别在第一、第三象限．  

(2) 若 $\alpha$、$\beta$ 分别在第二、第四象限，固定 $\alpha=\dfrac{3\pi}{4}$，于是等式变成了 $-1+\tan\beta=-\tan\beta(1+\tan\beta)$，即  
$$
\tan^2\beta+2\tan\beta-1=0.\quad (*)
$$  
考虑函数 $f(x)=x^2+2x-1$，由于 $f(-1)=-2$，$f(-10)>0$，所以 $f(x)$ 在区间 $(-10,-1)$ 中有零点，因此存在 $\beta$ 位于第四象限且使得 $(*)$ 成立．
\end{solutions}